% GENERATED FILE. Edit canonical lessons, then run npm run generate:books.
% canonical-source-hash: fnv1a64:ab14e760818d3e13
% canonical-lessons: ZH-C468-shanliang, ZH-C468-heqi, ZH-C468-jianqiang, ZH-C468-qinkuai, ZH-C468-titie

\chapter{Good Points}
\label{ch:zh-good-points}
\hlchaptermodality{468}

\begin{chapteropening}
\textbf{By the end of this chapter:} \emph{I can say someone is kind-hearted, good-natured, strong, hard-working or considerate.}

Complete the last lesson of chapter 468: I can say someone is kind-hearted, good-natured, strong, hard-working or considerate.
\end{chapteropening}

\section[shànliáng]{shànliáng — kind-hearted}

\label{lesson:ZH-C468-shanliang}



\begin{quote}
Before the new one: say the Mandarin for cheerful and outgoing, then the Mandarin for full of life.
\end{quote}

\subsection*{You'll want to know: shànliáng}
\emph{shànliáng} — \textquotedblleft{}kind-hearted\textquotedblright{}.

\begin{quote}
\emph{wǒ māma shì wǒ rènshi de zuì shànliáng de rén} — My mother is the kindest person I know.
\end{quote}

Its characters come later; for now it is a word for the ear.

\begin{cousinweb}
\emph{shànliáng} is built from \emph{shàn} and \emph{liáng}, both inherited from Middle Chinese, \emph{dzyenX} and \emph{ljang} (Guangyun, in Baxter's transcription), which glosses each with the other: 'good'.
\end{cousinweb}

\begin{grammarlens}[title={What you've built}]
One more word for saying how it was.
\end{grammarlens}

\subsection*{Your turn}
\begin{itemize}
\raggedright
  \item \emph{Say it:} \emph{shànliáng}
  \item \emph{Say it:} \emph{shànliáng}, once more
  \item \emph{Say it:} \emph{wǒ māma shì wǒ rènshi de zuì shànliáng de rén} — My mother is the kindest person I know.
  \item \emph{Recall:} say the Mandarin for cheerful and outgoing, then the Mandarin for full of life, then say \emph{shànliáng} again
\end{itemize}

\begin{tcolorbox}[breakable,colback=teal!4,colframe=teal!35!black,title={Before you move on}]
What does shànliáng mean? (\textquotedblleft{}Kind-hearted\textquotedblright{}.) Say it once more.
\end{tcolorbox}
\section[\texorpdfstring{\zh{和气} (héqi)}{和气 (héqi)}]{\zh{和气} (héqi) — good-natured; pleasant to deal with}

\label{lesson:ZH-C468-heqi}



\begin{quote}
Before the new one: say the Mandarin for full of life, then the Mandarin for kind-hearted.
\end{quote}

\subsection*{You'll want to know: \zh{和气}}
\textbf{\zh{和气}} — \emph{héqi} — \textquotedblleft{}good-natured; pleasant to deal with\textquotedblright{}.

\begin{quote}
\textbf{\zh{他很安静对人也很和气}} \emph{tā hěn ānjìng duì rén yě hěn héqi} — He's quiet, and he's pleasant to everyone too.
\end{quote}

\begin{cousinweb}
\textbf{\zh{和气}} (\emph{héqi}) is built from \zh{和} 'harmony' and \zh{气} 'breath', both inherited from Middle Chinese, \emph{hwa} and \emph{khj+jH} (Guangyun, in Baxter's transcription), which glosses the first 'in tune' and 'neither hard nor soft'. The second is said light.
\end{cousinweb}

\begin{grammarlens}[title={What you've built}]
One more word for saying how it was.
\end{grammarlens}

\subsection*{Your turn}
\begin{itemize}
\raggedright
  \item \emph{Say it:} \emph{héqi}
  \item \emph{Say it:} \emph{héqi}, once more
  \item \emph{Say it:} \emph{tā hěn ānjìng duì rén yě hěn héqi} — He's quiet, and he's pleasant to everyone too.
  \item \emph{Recall:} say the Mandarin for full of life, then the Mandarin for kind-hearted, then say \emph{héqi} again
\end{itemize}

\begin{tcolorbox}[breakable,colback=teal!4,colframe=teal!35!black,title={Before you move on}]
What does \zh{和气} mean? (\textquotedblleft{}Good-natured; pleasant to deal with\textquotedblright{}.) Say it once more.
\end{tcolorbox}
\section[jiānqiáng]{jiānqiáng — strong; tough (in spirit)}

\label{lesson:ZH-C468-jianqiang}



\begin{quote}
Before the new one: say the Mandarin for kind-hearted, then the Mandarin for pleasant to deal with.
\end{quote}

\subsection*{You'll want to know: jiānqiáng}
\emph{jiānqiáng} — \textquotedblleft{}strong; tough (in spirit)\textquotedblright{}.

\begin{quote}
\emph{tā bìng le liǎng nián kěshì tā hěn jiānqiáng} — She has been ill for two years, but she's very strong.
\end{quote}

Its characters come later; for now it is a word for the ear.

\begin{cousinweb}
\emph{jiānqiáng} is built from \emph{jiān} 'firm' and \emph{qiáng} 'strong', both inherited from Middle Chinese, \emph{ken} and \emph{gjang} (Guangyun, in Baxter's transcription), which glosses the first 'firm' and the second 'healthy' and 'strong'.
\end{cousinweb}

\begin{grammarlens}[title={What you've built}]
One more word for saying how it was.
\end{grammarlens}

\subsection*{Your turn}
\begin{itemize}
\raggedright
  \item \emph{Say it:} \emph{jiānqiáng}
  \item \emph{Say it:} \emph{jiānqiáng}, once more
  \item \emph{Say it:} \emph{tā bìng le liǎng nián kěshì tā hěn jiānqiáng} — She has been ill for two years, but she's very strong.
  \item \emph{Recall:} say the Mandarin for kind-hearted, then the Mandarin for pleasant to deal with, then say \emph{jiānqiáng} again
\end{itemize}

\begin{tcolorbox}[breakable,colback=teal!4,colframe=teal!35!black,title={Before you move on}]
What does jiānqiáng mean? (\textquotedblleft{}Strong; tough (in spirit)\textquotedblright{}.) Say it once more.
\end{tcolorbox}
\section[qínkuài]{qínkuài — diligent with chores; quick to pitch in}

\label{lesson:ZH-C468-qinkuai}



\begin{quote}
Before the new one: say the Mandarin for pleasant to deal with, then the Mandarin for strong in spirit.
\end{quote}

\subsection*{You'll want to know: qínkuài}
\emph{qínkuài} — \textquotedblleft{}diligent with chores; quick to pitch in\textquotedblright{}.

\begin{quote}
\emph{tā bǐ tā gēge qínkuài} — He's more hard-working than his big brother.
\end{quote}

Its characters come later; for now it is a word for the ear.

\begin{cousinweb}
\emph{qínkuài} is built from \emph{qín} 'toil' and \zh{快} 'quick', both inherited from Middle Chinese, \emph{gj+n} and \emph{khwaejH} (Guangyun, in Baxter's transcription), which glosses the first 'toil' and the second as what pleases the heart; the sense 'quick' for the second is not given there.
\end{cousinweb}

\begin{grammarlens}[title={What you've built}]
One more word for saying how it was.
\end{grammarlens}

\subsection*{Your turn}
\begin{itemize}
\raggedright
  \item \emph{Say it:} \emph{qínkuài}
  \item \emph{Say it:} \emph{qínkuài}, once more
  \item \emph{Say it:} \emph{tā bǐ tā gēge qínkuài} — He's more hard-working than his big brother.
  \item \emph{Recall:} say the Mandarin for pleasant to deal with, then the Mandarin for strong in spirit, then say \emph{qínkuài} again
\end{itemize}

\begin{tcolorbox}[breakable,colback=teal!4,colframe=teal!35!black,title={Before you move on}]
What does qínkuài mean? (\textquotedblleft{}Diligent with chores; quick to pitch in\textquotedblright{}.) Say it once more.
\end{tcolorbox}
\section[tǐtiē]{tǐtiē — considerate; thoughtful of others' needs}

\label{lesson:ZH-C468-titie}



\begin{quote}
Before the new one: say the Mandarin for strong in spirit, then the Mandarin for quick to pitch in.
\end{quote}

\subsection*{You'll want to know: tǐtiē}
\emph{tǐtiē} — \textquotedblleft{}considerate; thoughtful of others' needs\textquotedblright{}.

\begin{quote}
\emph{wǒ bìng de shíhou tā duì wǒ hěn tǐtiē} — When I was ill, he was very considerate to me.
\end{quote}

Its characters come later; for now it is a word for the ear.

\begin{cousinweb}
\emph{tǐtiē} is built from \zh{体} 'body' and \emph{tiē} 'to stick to', both inherited from Middle Chinese, \emph{thejX} and \emph{thep} (Guangyun, in Baxter's transcription). The dictionary glosses the first 'the body' and the second as pawning a thing for money; the sense 'considerate' is not given there.
\end{cousinweb}

\begin{grammarlens}[title={What you've built}]
One more word for saying how it was.
\end{grammarlens}

\subsection*{Your turn}
\begin{itemize}
\raggedright
  \item \emph{Say it:} \emph{tǐtiē}
  \item \emph{Say it:} \emph{tǐtiē}, once more
  \item \emph{Say it:} \emph{wǒ bìng de shíhou tā duì wǒ hěn tǐtiē} — When I was ill, he was very considerate to me.
  \item \emph{Recall:} say the Mandarin for strong in spirit, then the Mandarin for quick to pitch in, then say \emph{tǐtiē} again
\end{itemize}

\begin{tcolorbox}[breakable,colback=teal!4,colframe=teal!35!black,title={Before you move on}]
What does tǐtiē mean? (\textquotedblleft{}Considerate; thoughtful of others' needs\textquotedblright{}.) Say it once more.
\end{tcolorbox}
